\chapter{Construction d'indices et rééquilibrage}\label{ch:m1:index-construction-and-rebalancing}

À la clôture du 30 avril 2026, rien de visible ne se produisit~: la \omterm{def:m1:shares-corporate-actions-indices:share}{capitalisation
boursière} de chaque société américaine cotée fut consignée. Huit semaines plus tard, après
la clôture du 26 juin, les fonds qui suivent une même famille d'indices ont tous
négocié selon la liste produite par cet instantané, sur les mêmes titres, dans le
même sens, lors de la même enchère. Tout le monde connaissait la règle depuis un an et
la liste depuis un mois. Un indice est un algorithme publié, doté d'un calendrier, et
une somme d'argent considérable s'est engagée à s'y conformer. Ce chapitre lit
l'algorithme comme le ferait un trader~: ce qu'il faudra acheter, par qui,
quand, et quelle part de cette conséquence est déjà intégrée au cours.

\section{Un indice est un règlement de marché}

\begin{definition}[Méthodologie d'indice]\label{def:m1:index-construction-and-rebalancing:methodology}
Une \emph{méthodologie d'indice}\index{méthodologie d'indice} (index methodology) est le document publié
qui détermine un indice~: l'univers des titres éligibles, la
règle de sélection, la règle de pondération, le calendrier des révisions et le traitement
des opérations sur titres. La sélection est \emph{fondée sur des règles} lorsque le document
seul détermine les membres, et \emph{pilotée par un comité} lorsqu'un comité choisit
parmi les titres éligibles.
\end{definition}

\begin{definition}[Ajustement au flottant]\label{def:m1:index-construction-and-rebalancing:floatadj}
L'\emph{ajustement au flottant}\index{ajustement au flottant} (float adjustment) pondère chaque membre par sa
capitalisation multipliée par un facteur de flottant, soit la fraction de ses
actions dans le flottant (\cref{def:m1:shares-corporate-actions-indices:float}), de sorte
que l'indice puisse être détenu par tous à la fois.
\end{definition}

\begin{definition}[Reconstitution et rééquilibrage]\label{def:m1:index-construction-and-rebalancing:recon}
Une \emph{reconstitution d'indice}\index{reconstitution d'indice} (index reconstitution) est une révision programmée
au cours de laquelle la composition est de nouveau déterminée. Un \emph{rééquilibrage} met à jour
les pondérations (nombre d'actions, facteurs de flottant, plafonds) d'une composition inchangée.
Les deux suivent une séquence fixe~: un \emph{jour de classement} où les données sont
figées, une \emph{annonce} et une \emph{date d'effet}, après
la clôture de laquelle le nouvel indice s'applique.
\end{definition}

\begin{definition}[Indice pro forma]\label{def:m1:index-construction-and-rebalancing:proforma}
Un \emph{indice pro forma}\index{indice pro forma} (pro-forma index) est l'indice tel qu'il sera
après une révision à venir, composition et nombre d'actions, diffusé par le
fournisseur à ses clients entre l'annonce et la date d'effet, afin que les
trackers puissent préparer leurs transactions.
\end{definition}

\begin{dated}{2026-09}{Deux façons de retenir cinq cents ou deux mille actions}\label{dat:m1:index-construction-and-rebalancing:rules}
\textbf{Fondé sur des règles.} Les indices Russell US classent les sociétés par
capitalisation totale à la date de classement. Depuis 2026, ils se reconstituent deux fois par an.
Juin 2026~: date de classement le 30 avril, listes préliminaires mises à jour le 29 mai et les 5, 12 et 18
juin, effet après la clôture du vendredi 26 juin~; le point de rupture entre
le Russell 1000, à grande capitalisation, et le Russell 2000, à petite capitalisation,
était de 5.7~milliards de dollars, soit 24~\% de plus qu'un an plus tôt. En décembre 2026, la
première reconstitution de décembre depuis plus de trois décennies~: date de classement le 30 octobre,
listes préliminaires le 13 novembre, effet après la clôture du 11 décembre.
Le fournisseur évalue à environ 12.2~billions de dollars les actifs indexés sur ces
indices ou investis dans des produits qui s'y réfèrent.

\textbf{Piloté par un comité.} Les membres du S\&P 500 sont choisis par un comité
parmi les sociétés qui remplissent des critères publiés. Depuis le 1\up{er}~juillet 2025, une entrée
exige une capitalisation totale d'au moins 22.7~milliards de dollars (la ligne directrice est
révisée chaque trimestre et fixée près du 85\up{e}~centile de la \omterm{def:m1:shares-corporate-actions-indices:share}{capitalisation boursière}
cumulée), une capitalisation ajustée au flottant d'au moins la moitié de ce
seuil, et un ratio entre la valeur annuelle échangée et la capitalisation ajustée au
flottant d'au moins 0.75, outre des tests de rentabilité et
de domiciliation. Les critères s'appliquent aux entrées, et non au maintien dans l'indice.
\end{dated}

\begin{omfigure}
\begin{tikzpicture}[font=\small]
\draw[thick,-{Stealth[length=2.5mm]}] (0,0) -- (13,0);
\foreach \x/\d/\t in {0.8/30 avril/date de classement,5.0/29 mai/liste mise à jour,10.9/26 juin/{effet après la clôture}}{
  \draw[thick] (\x,0.12) -- (\x,-0.12);
  \node[above=3pt,font=\footnotesize\bfseries] at (\x,0.12) {\d};
  \node[below=3pt,font=\footnotesize,align=center] at (\x,-0.12) {\t};}
\foreach \x in {6.2,7.4,8.6,9.6}{\draw (\x,0.08) -- (\x,-0.08);}
\node[above=16pt,font=\footnotesize] at (7.9,0.12) {mises à jour ultérieures~: 5, 12, 18 juin};
\draw[decorate,decoration={brace,amplitude=5pt,mirror},omDef,thick] (0.8,-1.0) -- (4.9,-1.0) node[midway,below=6pt,font=\footnotesize,align=center,text=omDef] {chacun peut calculer la liste~:\\prévision};
\draw[decorate,decoration={brace,amplitude=5pt,mirror},omThm,thick] (5.1,-1.0) -- (10.9,-1.0) node[midway,below=6pt,font=\footnotesize,align=center,text=omThm] {la liste est publique~:\\positionnement, puis l'enchère de clôture};
\end{tikzpicture}

\omcaption{La reconstitution de juin 2026 d'une famille d'indices fondée sur des règles
(dates de l'encadré daté). Les cours postérieurs à la date de classement ne modifient pas la
composition~; ils modifient seulement le volume que les trackers devront
acheter.}\label{fig:m1:index-construction-and-rebalancing:timeline}
\end{omfigure}

\section{Règles de tampon}

Une règle naïve «~top $n$ par capitalisation~» fait entrer et sortir les titres proches
du rang $n$ à chaque révision. Chaque changement coûte aux trackers une transaction, donc
les fournisseurs amortissent la règle.

\begin{definition}[Règle de tampon]\label{def:m1:index-construction-and-rebalancing:buffer}
Une \emph{règle de tampon}\index{règle de tampon} (buffer rule) accorde aux membres en place un avantage à la
frontière~: avec un tampon $b$, un membre actuel d'un indice top-$n$ demeure aussi longtemps
qu'il se classe au rang $n+b$ ou mieux, et un non-membre n'entre que lorsqu'il se classe au
rang $n-b$ ou mieux~; les places restantes sont attribuées par rang. Les fournisseurs
expriment $b$ en rangs, en centiles de capitalisation cumulée, ou sous forme d'une bande
de pourcentage autour du point de rupture.
\end{definition}

\begin{omfigure}
\begin{tikzpicture}
\begin{axis}[width=11cm,height=6.2cm,
  xlabel={rang lors de cette révision},ylabel={rang lors de la révision précédente},
  tick label style={font=\small},label style={font=\small},
  xmin=148,xmax=252,ymin=80,ymax=420,grid=major,grid style={gray!25},
  legend columns=3,legend cell align=left,legend style={font=\footnotesize,at={(0.5,-0.27)},anchor=north,draw=none,fill=none,/tikz/every even column/.append style={column sep=6pt}},
  scatter/classes={stays={mark=*,omExo!70,mark size=1.4pt},outside={mark=o,omExo!70,mark size=1.4pt},
    added={mark=triangle*,omDef,mark size=2.6pt},deleted={mark=square*,omThm,mark size=2pt},
    kept={mark=*,omThm,mark size=2pt},blocked={mark=o,omDef,thick,mark size=2pt}},
  table/col sep=comma]
\addplot[scatter,only marks,scatter src=explicit symbolic] table[x=new_rank,y=old_rank,meta=outcome]{figdata/markets-1/15-index-construction-and-rebalancing/band.csv};
\draw[thick,dashed] (axis cs:200,80) -- (axis cs:200,420);
\draw[thin] (axis cs:180,80) -- (axis cs:180,420);
\draw[thin] (axis cs:220,80) -- (axis cs:220,420);
\legend{membre maintenu,hors indice,ajouté,retiré,membre conservé par le tampon,candidat bloqué par le tampon}
\end{axis}
\end{tikzpicture}

\omcaption{Une révision simulée d'un indice top-200 avec un tampon de 20
rangs (traits fins), pour les valeurs classées désormais entre les rangs 150 et 250. Entre les rangs 180
et 220, l'histoire décide~: les membres restent, les candidats attendent. Données~: la
simulation du tutoriel.}\label{fig:m1:index-construction-and-rebalancing:band}
\end{omfigure}

\begin{omfigure}
\begin{tikzpicture}
\begin{axis}[width=9.5cm,height=4.6cm,
  xlabel={tampon (rangs)},ylabel={ajouts par révision},
  tick label style={font=\small},label style={font=\small},
  xmin=-3,xmax=63,ymin=0,ymax=24,xtick={0,10,20,30,40,60},grid=major,grid style={gray!25},
  nodes near coords,nodes near coords style={font=\footnotesize,above=2pt},
  point meta=explicit,table/col sep=comma]
\addplot[omDef,very thick,mark=*] table[x=buffer,y=names,meta expr=\thisrow{weight_pct}]{figdata/markets-1/15-index-construction-and-rebalancing/turnover.csv};
\end{axis}
\end{tikzpicture}

\omcaption{Nombre moyen d'ajouts par révision annuelle d'un indice simulé
top-200 issu d'un univers de 600 valeurs, en fonction de la largeur du tampon. Le
nombre au-dessus de chaque point est le taux de rotation à sens unique, en pourcentage du poids de l'indice. Un
tampon de 30 rangs divise par deux le nombre de changements. Données~: la
simulation du tutoriel, 20 trajectoires de 20 ans.}\label{fig:m1:index-construction-and-rebalancing:turnover}
\end{omfigure}

\begin{method}[Prévoir une reconstitution fondée sur des règles]\label{met:m1:index-construction-and-rebalancing:predict}
\begin{enumerate}
\item \textbf{Univers.} Reconstruire l'univers éligible à partir de la méthodologie~:
place de cotation, domiciliation, cours minimum et flottant, catégories d'actions.
\item \textbf{Capitalisation.} Calculer la capitalisation totale comme le fait le fournisseur,
en tenant compte de toutes les catégories d'actions et des derniers nombres d'actions tirés des dépôts. La plupart
des erreurs de prévision sont des erreurs sur le nombre d'actions.
\item \textbf{Classement et tampon.} Classer les valeurs, repérer les seuils, appliquer le tampon
aux membres en place.
\item \textbf{Avant le jour du classement}, attribuer à chaque valeur frontière une probabilité
fonction de son écart au seuil et de sa volatilité sur les jours qui
restent.
\item \textbf{Dimensionner l'ordre} avec la \cref{prop:m1:index-construction-and-rebalancing:demand}.
\end{enumerate}
\end{method}

\section{Ce que doivent faire les trackers}

\begin{proposition}[Demande passive]\label{prop:m1:index-construction-and-rebalancing:demand}
Soient des fonds d'actifs $A$ répliquant un indice ajusté au flottant de capitalisation totale
ajustée au flottant $M$. Ils détiennent la même fraction $\varphi =
A/M$ du flottant de chaque membre. Lorsqu'une valeur de capitalisation ajustée au flottant
$F$ est ajoutée, ils doivent acheter $\varphi F$ dollars de ce titre, soit
une fraction $\varphi$ de son flottant (plus précisément $AF/(M+F)$, la différence
étant négligeable pour une seule addition).
\end{proposition}
\begin{proof}
Un tracker détient chaque membre à son poids dans l'indice $F_i/M$, soit $AF_i/M$
dollars du membre $i$, ce qui représente la fraction $A/M$ de $F_i$.
\end{proof}

\begin{definition}[Choc de demande]\label{def:m1:index-construction-and-rebalancing:shock}
Un \emph{choc de demande}\index{choc de demande} est une modification de la quantité
d'un titre que certains investisseurs doivent détenir pour des raisons sans lien avec sa valeur. Un
changement d'indice en est l'exemple le plus net~: la quantité, la date et l'identité
des acheteurs sont connus, et l'indice de référence des acheteurs est le cours de clôture de
la date d'effet, de sorte qu'ils sont indifférents au prix payé à cette
clôture.
\end{definition}

\begin{example}[Treize jours de volume]\label{ex:m1:index-construction-and-rebalancing:thirteen}
Un indice a $M = 5$ billions de dollars, et les trackers disposent de $A = 2$ billions de dollars, soit
$\varphi = 40\%$. Une société compte 400 millions d'actions cotées à \$50, dont 65~\% constituent le
flottant, soit $F = \$13$ milliards~: un poids de 0.26~\% et une demande de 5.2
milliards de dollars. Il s'échange pour 400 millions de dollars par jour~: les trackers
ont besoin de treize jours de volume moyen, en une seule clôture.
\end{example}

\begin{proposition}[L'opération de financement]\label{prop:m1:index-construction-and-rebalancing:funding}
Lorsqu'un ajout prend le poids $w$ dans l'indice, le poids de chaque membre restant
est multiplié par $1-w$~: les trackers vendent $Aw_iw$ du membre $i$, soit $Aw$
au total, moins la valeur de ce qui est retiré.
\end{proposition}

\begin{example}[Un grand ajout]\label{ex:m1:index-construction-and-rebalancing:large}
Le 16 novembre 2020, un fournisseur d'indices a annoncé qu'une société capable de peser
plus de 1~\% rejoindrait son indice phare de 500 valeurs. Il
a consulté le marché sur l'opportunité de l'ajouter en une ou deux étapes, puis le 30
novembre a annoncé une seule étape, au poids plein ajusté au flottant, en vigueur
avant l'ouverture du lundi 21 décembre, les fichiers pro forma étant diffusés
après la clôture du 11 décembre. Avec $A = 10$ billions de dollars et $w = 1.5\%$
(chiffres ronds, non ceux du fournisseur), les trackers achèteraient 150 milliards de dollars
d'une seule valeur et en vendraient 9 milliards de dollars d'un membre pesant 6~\%, le tout à une seule clôture~:
celle du vendredi 18 décembre, troisième vendredi du mois, et donc une échéance trimestrielle
des contrats à terme et options sur indices, parmi les clôtures les plus liquides de
l'année (voir au \cref{ch:m1:auctions}).
\end{example}

\section{L'effet indice, et ce qu'il est devenu}

\begin{definition}[Effet indice]\label{def:m1:index-construction-and-rebalancing:effect}
L'\emph{effet indice}\index{effet indice} est le rendement anormal
d'un titre entre l'annonce de son ajout à un indice (ou de son retrait de celui-ci) et
la date d'effet, ainsi que son retournement partiel ensuite.
\end{definition}

Si les investisseurs avaient une courbe de demande plate pour chaque titre, un \omterm{def:m1:index-construction-and-rebalancing:shock}{choc de demande}
sans contenu informatif ne ferait pas bouger le cours. Pendant des décennies, c'était le cas~: une étude
des ajouts au S\&P 500 trouve un rendement anormal moyen de 7.4~\% dans les
années 1990. La même étude trouve 0.3~\% pour 2010--2020, et pour les retraits un fort
rendement négatif dans les années 1990, contre 0.1~\% pour la décennie suivante, alors que
les actifs répliquant l'indice s'étaient multipliés dans l'intervalle. Le titre de l'étude
est sa conclusion~: l'effet a disparu de la fenêtre d'annonce.

Trois explications, non exclusives. \emph{Prévision}~: les ajouts sont désormais anticipés, si bien
que le cours s'ajuste avant l'annonce, là où une étude d'événement
ne regarde pas. \emph{Migration}~: beaucoup d'ajouts à l'indice large
proviennent de l'indice des valeurs moyennes du fournisseur, dont les trackers vendent ce que
les trackers de l'indice large achètent. \emph{Offre de liquidité}~: une population
d'entreprises existe désormais pour acheter en amont et vendre aux trackers à la clôture. Le
\omterm{def:m1:index-construction-and-rebalancing:shock}{choc de demande} n'a pas disparu~; son prix a été laminé par la concurrence, et
il s'est déplacé plus tôt.

\begin{omfigure}
\begin{tikzpicture}
\begin{axis}[width=11cm,height=5cm,
  xlabel={jours de négociation depuis la date d'effet},ylabel={rendement anormal (\%)},
  tick label style={font=\small},label style={font=\small},
  xmin=-15,xmax=30,ymin=-0.5,ymax=9,grid=major,grid style={gray!25},
  legend columns=2,legend style={font=\footnotesize,at={(0.5,-0.3)},anchor=north,draw=none,fill=none,/tikz/every even column/.append style={column sep=8pt}},
  table/col sep=comma]
\addplot[omThm,very thick] table[x=day,y=then_pct]{figdata/markets-1/15-index-construction-and-rebalancing/event.csv};
\addplot[omDef,very thick,dashed] table[x=day,y=now_pct]{figdata/markets-1/15-index-construction-and-rebalancing/event.csv};
\legend{ramené à la moyenne des années 1990 (7.4~\%),ramené à la moyenne 2010--2020 (0.3~\%)}
\end{axis}
\end{tikzpicture}

\omcaption{Un ajout stylisé, en rendement anormal cumulé~: un saut à l'annonce (jour $-5$), une
hausse progressive jusqu'à la clôture effective, un retournement partiel. Seules les deux hauteurs au
jour 0 sont des moyennes publiées~; les formes, la répartition entre parties permanente et
temporaire et la vitesse du retournement sont
illustratives.}\label{fig:m1:index-construction-and-rebalancing:event}
\end{omfigure}

Les événements d'indice restent une activité lucrative (le livre~8 fournit les fichiers de stratégie)
parce que les moyennes masquent une large dispersion~: des reconstitutions de petites capitalisations
dont les besoins atteignent plusieurs jours de volume, des migrations, des changements de flottant et de
nombre d'actions que personne n'annonce dans les gros titres, et des indices sur des
marchés où moins d'entreprises font le calcul.

\section{Tutoriel~: une reconstitution avec tampons}\label{tut:m1:index-construction-and-rebalancing:recon}

\begin{tutorial}
\textbf{Objectif.} Reconstituer un indice des $n$ premières valeurs avec un tampon, mesurer ce que le
tampon fait économiser et dimensionner la demande des trackers. \textbf{État final~:} les trois
résultats chiffrés de ce chapitre.
\begin{tutsteps}
\item \textbf{La règle.} Conserver les membres dans le rang $n+b$, admettre les entrantes dans le rang
$n-b$, puis compléter ou réduire jusqu'à exactement $n$.
=\omcode{code/markets-1/15-index-construction-and-rebalancing/python/index_recon.py}{15}{31}{Une reconstitution des $n$ premières valeurs avec un tampon de $b$ rangs.}\label{lst:m1:index-construction-and-rebalancing:recon}
\item \textbf{Rotation.} Simuler 600 capitalisations lognormales avec une volatilité annuelle de 35~\%
et une révision une fois par an
(\cref{fig:m1:index-construction-and-rebalancing:turnover}).
\item \textbf{Demande.}
=\omcode{code/markets-1/15-index-construction-and-rebalancing/python/index_recon.py}{53}{59}{Ce que les trackers doivent acheter, en dollars et en jours de volume.}\label{lst:m1:index-construction-and-rebalancing:demand}
\end{tutsteps}
\textbf{Pour aller plus loin.} Remplacer le tampon de rang par une bande de $\pm
2.5\%$ de capitalisation cumulée autour du point de rupture, puis comparer la
rotation. Donner ensuite à chaque valeur frontière une probabilité de franchissement entre
aujourd'hui et le jour de rang, à partir de sa volatilité, et vérifier la calibration de
vos probabilités sur les années simulées.
\end{tutorial}

\section{Construction~: le prédicteur de reconstitution}\label{bld:m1:index-construction-and-rebalancing:predictor}

\begin{build}
\bfield{Rôle} Le carnet d'événements de l'entreprise miniature (le livre~8) part
d'une liste~: quels titres entreront et sortiront, et combien de jours de volume
les trackers devront traiter pour chacun. \bfield{Interface} \texttt{Security(symbol, shares, price, float\_factor,
adv\_shares, member)}~; \texttt{predict(universe, n, buffer)} renvoyant
les membres, les ajouts et les retraits~; \texttt{weights(universe, members)},
pondérations ajustées au flottant~; \texttt{tracker\_trades(universe, prediction,
tracked\_assets)} renvoyant pour chaque ajout et chaque retrait les dollars, les
actions et le multiple du volume quotidien moyen.
\bfield{Règles} Classement sur la capitalisation totale, pondération sur la capitalisation
ajustée au flottant~: les deux sont différentes à dessein. Les ex aequo sont départagés par symbole afin
que la sortie soit déterministe. Les nombres d'actions proviennent de l'ajusteur d'opérations sur
titres (\cref{bld:m1:shares-corporate-actions-indices:adjuster}).
\bfield{Tests d'acceptation} \texttt{code/firm/recon/tests/}~: un univers de cinq valeurs
établi à la main, avec et sans tampon~; pondérations ajustées au flottant~; les
transactions des trackers lors d'un ajout et d'un retrait.
\bfield{Pour aller plus loin} Probabilités de franchissement avant le jour de rang~; migrations
entre deux indices de la même famille, avec des actifs suivis différents.
\end{build}

\begin{omsources}
\item FTSE Russell, \emph{FTSE Russell begins June 2026 semi-annual Russell
US Indexes reconstitution} et \emph{FTSE Russell announces December 2026
Russell US Indexes reconstitution schedule}, communiqués de presse, 2026.
\item S\&P Dow Jones Indices, \emph{Update to S\&P Composite 1500 Market Cap
Guidelines}, communiqués de presse des 4 janvier 2023 et 1\up{er}~juillet 2025~; \emph{S\&P Dow
Jones Indices Announces Implementation of Tesla's Addition to S\&P 500}, 30
novembre 2020.
\item R.~Greenwood et M.~Sammon, ``The disappearing index effect'',
\emph{Journal of Finance} 80 (2025), 657--698.
\end{omsources}

\section{Exercices}

\begin{exercise}[$\star$]\label{exo:m1:index-construction-and-rebalancing:1}
Une entreprise a 400 millions d'actions à 50~\$~; son fondateur en détient 35~\% et le reste
est dans le flottant. Calculer sa capitalisation ajustée au flottant et son poids dans un indice
dont la capitalisation ajustée au flottant est de 5~000 milliards de dollars.
\end{exercise}

\begin{exercise}[$\star$]\label{exo:m1:index-construction-and-rebalancing:2}
Des fonds de 2~000 milliards de dollars répliquent cet indice. Quelle fraction du flottant de
chaque membre détiennent-ils~? L'entreprise de l'exercice précédent est ajoutée et s'échange pour
400 millions de dollars par jour. Calculer la demande en dollars et en jours de volume.
\end{exercise}

\begin{exercise}[$\star$]\label{exo:m1:index-construction-and-rebalancing:3}
Un indice des 100 premières valeurs a un tampon de 10 rangs. Décider du sort de~: un membre
classé 108e~; une entrante classée 95e~; une entrante classée 88e~; un membre classé 112e.
\end{exercise}

\begin{exercise}[$\star\star$]\label{exo:m1:index-construction-and-rebalancing:4}
Lors de la reconstitution de juin 2026 de
l'\cref{dat:m1:index-construction-and-rebalancing:rules}, une entreprise est classée
juste au-dessus du point de rupture de 5.7~milliards de dollars au jour de rang et perd 30~\% de
sa valeur en mai. Dans quel indice se trouve-t-elle le 29 juin~? Qu'a changé la baisse pour
les fonds qui répliquent cet indice, et pour une entreprise qui a acheté le titre en avril
en comptant sur leurs achats~?
\end{exercise}

\begin{exercise}[$\star\star$]\label{exo:m1:index-construction-and-rebalancing:5}
Avec les chiffres ronds de
l'\cref{ex:m1:index-construction-and-rebalancing:large}, vérifier les deux montants en dollars.
Quel est le nouveau poids du membre à 6~\%~? Pourquoi le fournisseur a-t-il pu envisager
deux étapes, et quel est l'argument pour une seule~?
\end{exercise}

\begin{exercise}[$\star\star$]\label{exo:m1:index-construction-and-rebalancing:6}
Appliquer la règle de la racine carrée, telle que présentée à la \cref{met:m1:the-sell-side:sqrt}, avec
$\kappa = 0.7$ et une volatilité quotidienne de 2.5~\% à la demande de treize
jours de volume de l'\cref{ex:m1:index-construction-and-rebalancing:thirteen}, comme
si elle était exécutée en un seul jour. Supposer ensuite que des firmes qui anticipent achètent la même
quantité, uniformément, sur les vingt jours précédents et la cèdent à la clôture~:
estimer l'impact de chacun de leurs jours. Comparer aux deux moyennes
citées dans ce chapitre.
\end{exercise}

\begin{exercise}[$\star\star\star$]\label{exo:m1:index-construction-and-rebalancing:7}
\emph{Programmation.} Avec \texttt{simulate\_turnover(600, 200, b, 20, 1)}, indiquer
les ajouts par revue et le taux de rotation unilatéral pour $b = 0$ et $b = 30$.
Un fonds indiciel paie 15~points de base sur ce qu'il négocie~; que fait
économiser la \omterm{def:m1:index-construction-and-rebalancing:buffer}{règle de tampon} à ses investisseurs par an, en points de base d'actifs~?
\end{exercise}

\begin{exercise}[$\star\star\star$]\label{exo:m1:index-construction-and-rebalancing:8}
\emph{Trouver la faille.} «~Acheter à l'ouverture suivante chaque ajout annoncé à l'indice des grandes
capitalisations et le vendre à la clôture effective a rapporté 5~\% par événement dans mon
backtest sur 1990--2000. Il y a environ vingt-cinq événements par an. Je m'attends à
125~\% par an, sans levier.~» Donner trois raisons indépendantes pour lesquelles
cette espérance est fausse.
\end{exercise}

\section{Problème~: Jour de reconstitution}

\begin{problem}[{Problème du week-end --- une petite valeur le dernier vendredi de juin}]\label{pb:m1:index-construction-and-rebalancing:1}
Un indice à petite capitalisation affiche une \omterm{def:m1:shares-corporate-actions-indices:share}{capitalisation boursière} ajustée au flottant
de 3~billions de dollars et des trackers pour 200~milliards de dollars. Un indice à grande capitalisation de la
même famille affiche 50~billions de dollars et des trackers pour 4~billions de dollars. Le titre X (cours
24~\$, capitalisation ajustée au flottant de 1.2 milliards de dollars, volume moyen de 0.9 million
d'actions par jour, volatilité quotidienne de 3~\%) doit rejoindre le petit indice. Le titre Y (cours
60~\$, flottant de 9~milliards de dollars, volume moyen de 2.5 millions d'actions) passe du
petit indice vers le grand.

\medskip\noindent\textbf{Partie I --- La demande.}
\begin{enumerate}
\item Quelle fraction du flottant de chaque membre détiennent les trackers du petit
indice~? Et ceux du grand indice~?
\item Donner la demande des trackers pour X en dollars et en actions.
\item L'exprimer en jours de volume moyen.
\item Pour Y, donner la vente des trackers du petit indice, l'achat des trackers du grand et
le solde.
\item Exprimer la vente et le solde en jours de volume de Y. Quel nombre
compte pour le cours de clôture~?
\end{enumerate}

\medskip\noindent\textbf{Partie II --- La clôture.}
\begin{enumerate}[resume]
\item Les trackers envoient 60~\% de leur demande pour X à l'enchère de clôture. Chez X,
l'enchère de clôture porte normalement sur 8~\% du volume d'une journée. Comparer.
\item Estimer l'impact de la demande totale pour X avec la règle de la racine
carrée, $\kappa = 0.7$, comme si elle arrivait en une seule journée.
\item Même question si elle est répartie uniformément sur vingt jours.
\item Qui vend 2 millions d'actions de X à l'enchère, et d'où
les tient-il~?
\end{enumerate}

\medskip\noindent\textbf{Partie III --- La transaction d'anticipation.}
Trois semaines avant le jour de rang, une firme croit que X sera ajoutée avec
une probabilité de 0.8. Si elle est ajoutée, elle attend une hausse de 3~\% jusqu'à la clôture effective~;
sinon, une baisse de 2~\% lorsque les autres anticipateurs se retirent. Un aller-retour coûte 40
points de base. Elle achète pour 5~millions de dollars.
\begin{enumerate}[resume]
\item Donner le rendement espéré et le gain espéré.
\item Donner l'écart-type du gain.
\item La firme détient 25 positions de ce type. Les résultats étant indépendants, donner
le gain espéré, son écart-type et leur rapport.
\item Même question avec une corrélation deux à deux de 0.3 entre les rendements des positions.
\item À quelle probabilité d'ajout la transaction atteint-elle le seuil de rentabilité~?
\item Citer deux sources de cette corrélation.
\end{enumerate}

\medskip\noindent\textbf{Partie IV --- Appréciation.}
\begin{enumerate}[resume]
\item Le fournisseur passe d'une reconstitution par an à deux. Que se passe-t-il
pour la taille de chaque événement et pour l'activité des anticipateurs~?
\item Pourquoi les trackers acceptent-ils de payer l'impact à la clôture plutôt que
d'acheter plus tôt~?
\item Le gérant d'un fonds indiciel \emph{achète bien} en avance. Qui supporte le risque, et
comment ses investisseurs s'en apercevraient-ils~?
\item Énoncer le \emph{résultat nommé}~: la demande des trackers pour X, exprimée en multiple
de son volume quotidien moyen.
\item En une phrase~: quelle est la fonction économique des firmes qui
anticipent les changements d'indice~?
\end{enumerate}
\end{problem}

\section{Questions d'entretien}

\begin{interviewq}[$\star$ \iqroles{trader, chercheur}]\label{iq:m1:index-construction-and-rebalancing:1}
Un titre entre dans un grand indice. Que se passe-t-il pour son cours, et quand~?
\end{interviewq}

\begin{interviewq}[$\star$ \iqroles{chercheur, développeur}]\label{iq:m1:index-construction-and-rebalancing:2}
Pourquoi les indices sont-ils ajustés au flottant, et qu'est-ce qui casserait s'ils ne l'étaient pas~?
\end{interviewq}

\begin{interviewq}[$\star\star$ \iqroles{chercheur, trader}]\label{iq:m1:index-construction-and-rebalancing:3}
Les trackers détiennent 8~\% d'un indice. Un titre dont le flottant vaut 10~milliards de dollars est ajouté.
Combien doivent-ils acheter~? Que devez-vous savoir d'autre pour dire si cela est
beaucoup~?
\end{interviewq}

\begin{interviewq}[$\star\star$ \iqroles{chercheur}]\label{iq:m1:index-construction-and-rebalancing:4}
Comment prévoiriez-vous les ajouts de juin prochain dans un indice à petite capitalisation fondé sur des règles,
et d'où viendraient vos erreurs~?
\end{interviewq}

\begin{interviewq}[$\star\star$ \iqroles{chercheur, trader}]\label{iq:m1:index-construction-and-rebalancing:5}
Les encours passifs ont énormément augmenté, pourtant l'\omterm{def:m1:index-construction-and-rebalancing:effect}{effet indice} mesuré a
diminué. Concilier ces deux faits.
\end{interviewq}

\begin{interviewq}[$\star\star\star$ \iqroles{trader, chercheur}]\label{iq:m1:index-construction-and-rebalancing:6}
Vous dirigez un carnet de rééquilibrage d'indices. Décrivez ses principaux risques et comment vous
dimensionneriez les positions.
\end{interviewq}
